\section{The Hour We Have Not Been Told}

The certainty of death and the uncertainty of its hour belong together. We can assent to the first while living as though the second guaranteed delay. Yet an unknown date supplies no promise of more time. In the parable of the rich man, the barns are enlarged and the years of enjoyment confidently anticipated before God announces that the man's life will be required that night (Lk 12:15--21). The error lies in treating stored goods as security for a future the owner cannot command. His preparations have left the most necessary relation unattended: he is not rich toward God.

Prudent provision remains a duty. Dependents need to be supported; responsibilities may require saving, planning, and patient work. But we can make provision as though the continuation of our life were assured. We may consider every foreseeable interruption to our plans while neglecting our dependence on God for life itself. We can examine this danger by asking what our preparations serve: the needs of those entrusted to us, or a confidence in our arrangements that has allowed prayer and repentance to be put off.

St.~Francis de Sales considers the uncertain hour in the fifth meditation of the \emph{Introduction to the Devout Life} (I.13). He asks the reader to consider that death may come in one season or another, in one place or another, with warning or suddenly, with a confessor at hand or without that opportunity. Each possibility makes it harder to assume that we will recognize a safe time to begin preparing. A reader who now has access to confession and the ability to amend a fault can consider how much depends on circumstances that may later be withdrawn.

He then considers the farewells death will require: to possessions and pleasures, to companions and family, and finally to one's own body. He asks the reader to consider how these attachments will appear at death and to choose holy friendship, a love capable of enduring in God. This choice requires attention to the way we love now. Francis directs the reader to examine conscience, put what is disordered into order, and seek the help of the Blessed Virgin and the saints. His meditation ends in resolutions for the life the reader can still amend.

Alphonsus Liguori presses the same point in \emph{Preparation for Death}, especially considerations IV.1 and V.1--2. A person may know perfectly well that everyone dies and still place his own death at a convenient distance. Alphonsus warns especially against intending to sin now and confess tomorrow. Such a decision presumes a future opportunity to repent. It also presumes that we will later be willing to relinquish the sin we are deliberately choosing today. We can guarantee neither. The time God gives for repentance is being misused when we take it as permission to continue refusing him.

We can apply this warning to a decision already before us. We may be postponing confession because the sin is embarrassing, delaying restitution because it will cost money, or putting off a difficult reconciliation until the other person becomes easier to approach. Embarrassment, cost, and difficulty will not necessarily diminish with time. Continuing to avoid them can strengthen the habit of refusal. Acknowledging the cost of repentance allows us to ask for the grace to bear it and take the first step, while the opportunity is still available.

In Luke's account, the servants keep their lamps burning and wait for their master. When he finds them ready, he seats them at table and serves them (Lk 12:35--40). Christ commands watchfulness with this promise of welcome. The servants have reason to desire their master's arrival, because he will receive them with a generosity beyond the service they have rendered. We too can await the Lord with desire as we seek to be ready for him, asking that our duties be fulfilled in hope of that welcome.

An annual meditation can expose a pattern that a single earnest resolution misses. The same postponed confession, entrenched resentment, or neglected obligation may return before us year after year. Time has been passing while the excuse has remained almost unchanged. Recognition need not lead to despair. It gives a specific place to begin. Instead of making another broad promise to become serious, we can perform the next concrete act that seriousness already requires. A priest can be contacted, a debt acknowledged, an occasion of sin removed, or a neglected person visited.

We cannot recover the years already spent. We can seek mercy for our sins and use the time that remains to amend our life. Alphonsus addresses a reader who is still able to respond. If his warning makes us afraid, we can bring that fear to confession and prayer, asking Christ for help with the particular sin we have been unwilling to leave. The opportunity to repent is itself a reason for gratitude.

